\documentclass[12pt, a4paper]{article}
\usepackage[UTF8]{ctex}
\usepackage{geometry}
\usepackage{amsmath}
\usepackage{amssymb}
\usepackage{graphicx}
\usepackage{listings}
\usepackage{xcolor}
\usepackage{booktabs}
\usepackage{hyperref}
\usepackage{float}
\usepackage{subfigure}
\usepackage{fancyhdr}
\usepackage{setspace}

\geometry{left=2.5cm, right=2.5cm, top=3cm, bottom=3cm, headheight=14.5pt}

\definecolor{codebg}{rgb}{0.95,0.95,0.95}
\definecolor{codekey}{rgb}{0.0,0.4,0.8}
\definecolor{codecomment}{rgb}{0.3,0.6,0.3}
\definecolor{codestring}{rgb}{0.8,0.2,0.2}

\lstset{
    language=Python,
    backgroundcolor=\color{codebg},
    basicstyle=\small\ttfamily,
    keywordstyle=\color{codekey}\bfseries,
    commentstyle=\color{codecomment}\itshape,
    stringstyle=\color{codestring},
    showstringspaces=false,
    breaklines=true,
    frame=single,
    framerule=0.5pt,
    rulecolor=\color{gray},
    numbers=left,
    numberstyle=\tiny\color{gray},
    numbersep=8pt,
    tabsize=4,
    captionpos=b
}

\pagestyle{fancy}
\fancyhf{}
\rhead{深度学习 作业4}
\lhead{学号：2353224 \quad 姓名：魏佳怡}
\cfoot{\thepage}

\begin{document}

\begin{titlepage}
    \centering
    \vspace*{3cm}
    {\Huge\bfseries 深度学习实验报告}\\
    \vspace{0.5cm}
    {\Large 第七章：Seq2Seq 与注意力机制}\\
    \vspace{2cm}
    \begin{tabular}{ll}
        \large\textbf{学号：} & \large 2353224 \\\\
        \large\textbf{姓名：} & \large 魏佳怡 \\\\
        \large\textbf{课程：} & \large 深度学习 \\\\
        \large\textbf{实验题目：} & \large 序列到序列模型与注意力机制 \\\\
        \large\textbf{日期：} & \large 2026年3月29日 \\\\
    \end{tabular}
    \vfill
\end{titlepage}

\onehalfspacing

%=============================================================
\section{实验目的}
%=============================================================

本次实验旨在通过构建序列逆置任务，深入理解以下核心内容：

\begin{enumerate}
    \item 理解 Sequence-to-Sequence（Seq2Seq）模型的基本结构，包括 Encoder 和 Decoder 的工作原理；
    \item 掌握 Teacher Forcing 训练策略以及自回归推理（Autoregressive Decoding）的区别与联系；
    \item 理解注意力机制（Attention Mechanism）的动机、原理及其对 Seq2Seq 模型的改进；
    \item 实现双线性注意力（Bilinear Attention），并将其融入 Seq2Seq 框架；
    \item 对比两种模型在序列逆置任务上的训练效率与泛化能力。
\end{enumerate}

%=============================================================
\section{实验原理}
%=============================================================

\subsection{Sequence-to-Sequence 模型}

Seq2Seq 模型由 Sutskever 等人于2014年提出，是处理变长输入输出的经典框架。其核心思想是将输入序列通过一个 \textbf{Encoder} 压缩为固定维度的语义向量（context vector），再由 \textbf{Decoder} 根据该向量逐步生成目标序列。

\subsubsection{Encoder}

Encoder 使用循环神经网络（RNN）对输入序列 $\mathbf{x} = (x_1, x_2, \ldots, x_T)$ 逐步处理，在每个时间步更新隐状态：

\begin{equation}
    h_t^{\text{enc}} = f_{\text{enc}}(x_t, h_{t-1}^{\text{enc}})
\end{equation}

最终将最后一个时间步的隐状态 $h_T^{\text{enc}}$ 作为整个输入序列的语义表示，传递给 Decoder。

\subsubsection{Decoder}

Decoder 同样使用 RNN，以 Encoder 的最终隐状态初始化，并在每个时间步预测下一个输出 token：

\begin{equation}
    h_t^{\text{dec}} = f_{\text{dec}}(y_{t-1}, h_{t-1}^{\text{dec}})
\end{equation}

\begin{equation}
    P(y_t \mid y_{<t}, \mathbf{x}) = \text{softmax}(W h_t^{\text{dec}})
\end{equation}

\subsubsection{Teacher Forcing}

在训练阶段，Decoder 的输入使用真实标签序列（ground truth），而非上一步的预测结果，这种策略称为 \textbf{Teacher Forcing}。其优点是训练稳定、收敛快；缺点是训练与推理阶段存在分布偏移（exposure bias）。

在推理阶段，Decoder 采用自回归方式，将上一步预测的 token 作为当前步的输入，逐步生成完整输出序列。

\subsection{注意力机制}

传统 Seq2Seq 模型的瓶颈在于将整个输入序列压缩为单一固定维度向量，当序列较长时信息损失严重。Bahdanau 等人（2015）提出的注意力机制解决了这一问题：Decoder 在每个解码步骤可以"选择性地关注"Encoder 各时刻的输出，而不仅依赖最终隐状态。

\subsubsection{双线性注意力（Bilinear Attention）}

本实验采用双线性注意力（也称 Luong Attention 的一种变体），其注意力分数计算如下：

\begin{equation}
    \text{score}(h_t^{\text{dec}}, h_s^{\text{enc}}) = h_t^{\text{dec}} \cdot W \cdot (h_s^{\text{enc}})^\top
\end{equation}

其中 $W \in \mathbb{R}^{d \times d}$ 是可学习的投影矩阵（由 Dense 层实现）。

\subsubsection{注意力权重与上下文向量}

对所有 Encoder 时刻的分数进行 softmax 归一化，得到注意力权重：

\begin{equation}
    \alpha_{t,s} = \frac{\exp(\text{score}(h_t^{\text{dec}}, h_s^{\text{enc}}))}{\sum_{s'=1}^{T} \exp(\text{score}(h_t^{\text{dec}}, h_{s'}^{\text{enc}}))}
\end{equation}

再以注意力权重对 Encoder 输出做加权求和，得到上下文向量（context vector）：

\begin{equation}
    c_t = \sum_{s=1}^{T} \alpha_{t,s} \cdot h_s^{\text{enc}}
\end{equation}

\subsubsection{输出预测}

将 Decoder 当前隐状态 $h_t^{\text{dec}}$ 与上下文向量 $c_t$ 拼接后，通过线性层预测输出：

\begin{equation}
    \text{logits}_t = W_o \cdot [h_t^{\text{dec}};\, c_t]
\end{equation}

%=============================================================
\section{实验设置}
%=============================================================

\subsection{任务描述}

本实验以\textbf{字符串序列逆置}为玩具任务。输入为随机生成的大写字母字符串（仅含 A--Z），目标是输出其逆序。例如：

\begin{center}
    输入：\texttt{OIMESIQFIQ} \quad $\longrightarrow$ \quad 输出：\texttt{QIFQISEMIO}
\end{center}

字符采用 1--26 的整数编码（A=1, B=2, ..., Z=26），解码器起始 token 设为 0（特殊符号）。

\subsection{数据生成}

每个训练批次随机生成长度为 $L$ 的字符串，批量大小为 32。数据格式包含：

\begin{itemize}
    \item \texttt{enc\_x}：Encoder 输入，形状 $(B, L)$
    \item \texttt{dec\_x}：Decoder 输入（Teacher Forcing），形状 $(B, L)$，首位为起始符 0
    \item \texttt{y}：Decoder 目标输出（逆序序列），形状 $(B, L)$
\end{itemize}

\subsection{模型超参数}

\begin{table}[H]
    \centering
    \caption{两个模型的超参数设置}
    \begin{tabular}{lcc}
        \toprule
        超参数 & 普通 Seq2Seq & Attentive Seq2Seq \\
        \midrule
        词表大小 $|V|$ & 27 & 27 \\
        Embedding 维度 & 64 & 64 \\
        RNN 隐层维度 $d$ & 128 & 128 \\
        Attention 投影维度 & -- & 128 \\
        输出层输入维度 & 128 & 256（拼接后）\\
        训练步数 & 3000 & 2000 \\
        批量大小 & 32 & 32 \\
        训练序列长度 & 20 & 20 \\
        优化器 & Adam & Adam \\
        学习率 & 0.0005 & 0.0005 \\
        \bottomrule
    \end{tabular}
\end{table}

\subsection{实验环境}

\begin{itemize}
    \item 操作系统：Windows 10
    \item Python 版本：3.12.7
    \item 深度学习框架：TensorFlow / Keras 3
    \item 开发工具：Jupyter Notebook（Cursor IDE）
\end{itemize}

%=============================================================
\section{模型实现}
%=============================================================

\subsection{任务一：普通 Seq2Seq 模型}

\subsubsection{模型结构}

普通 Seq2Seq 模型由以下组件构成：

\begin{itemize}
    \item \texttt{Embedding}：将离散 token 映射为 64 维稠密向量
    \item \texttt{SimpleRNNCell}（Encoder）：隐层维度 128
    \item \texttt{SimpleRNNCell}（Decoder）：隐层维度 128
    \item \texttt{Dense}：将 Decoder 隐状态映射到词表大小（27维）的 logits
\end{itemize}

\subsubsection{前向传播实现}

\begin{lstlisting}[caption={普通 Seq2Seq 模型 call 方法}]
def call(self, enc_ids, dec_ids):
    # Encoder: 将输入序列编码
    enc_emb = self.embed_layer(enc_ids)  # (b_sz, enc_len, emb_sz)
    enc_out, enc_state = self.encoder(enc_emb)  # enc_state: (b_sz, h_sz)

    # Decoder: teacher forcing，用 encoder 最后隐状态初始化
    dec_emb = self.embed_layer(dec_ids)  # (b_sz, dec_len, emb_sz)
    dec_out, _ = self.decoder(dec_emb, initial_state=[enc_state])

    logits = self.dense(dec_out)  # (b_sz, dec_len, v_sz)
    return logits
\end{lstlisting}

Encoder 将输入序列处理完毕后，取最终隐状态 \texttt{enc\_state} 作为 Decoder 的初始状态，Decoder 在 Teacher Forcing 模式下逐步生成每个时间步的 logits。

\subsubsection{推理阶段单步解码}

推理时，Decoder 逐步生成 token，每步将上一步预测的 token 作为当前输入：

\begin{lstlisting}[caption={单步解码 get\_next\_token}]
def get_next_token(self, x, state):
    inp_emb = self.embed_layer(x)          # (b_sz, emb_sz)
    h, state = self.decoder_cell(inp_emb, state)  # (b_sz, h_sz)
    logits = self.dense(h)                 # (b_sz, v_sz)
    out = tf.argmax(logits, axis=-1)
    return out, state
\end{lstlisting}

\subsection{任务二：带注意力机制的 Seq2Seq 模型}

\subsubsection{模型结构}

在普通 Seq2Seq 基础上，增加以下组件：

\begin{itemize}
    \item \texttt{dense\_attn}：注意力投影层，实现双线性注意力中的矩阵 $W$，维度 $128 \to 128$
    \item \texttt{dense}：输出层输入维度变为 256（Decoder 隐状态 128 维 + context 向量 128 维拼接）
\end{itemize}

\subsubsection{前向传播实现}

\begin{lstlisting}[caption={Attentive Seq2Seq 模型 call 方法}]
def call(self, enc_ids, dec_ids):
    # --- Encoder ---
    enc_emb = self.embed_layer(enc_ids)           # (b, enc_len, emb_sz)
    enc_out, enc_state = self.encoder(enc_emb)    # enc_out: (b, enc_len, h)

    # --- Decoder (teacher forcing) ---
    dec_emb = self.embed_layer(dec_ids)           # (b, dec_len, emb_sz)
    dec_out, _ = self.decoder(dec_emb, initial_state=[enc_state])

    # --- Bilinear Attention ---
    enc_proj = self.dense_attn(enc_out)           # W * enc_out: (b, enc_len, h)
    scores = tf.matmul(dec_out, enc_proj, transpose_b=True)  # (b, dec_len, enc_len)
    attn_weights = tf.nn.softmax(scores, axis=-1) # (b, dec_len, enc_len)
    context = tf.matmul(attn_weights, enc_out)    # (b, dec_len, h)

    # 拼接后预测
    combined = tf.concat([dec_out, context], axis=-1)  # (b, dec_len, 2h)
    logits = self.dense(combined)                       # (b, dec_len, v_sz)
    return logits
\end{lstlisting}

\subsubsection{带注意力的单步解码}

推理阶段，每步解码需要额外传入 Encoder 的全部输出 \texttt{enc\_out}，以计算当前步的注意力：

\begin{lstlisting}[caption={带注意力的单步解码 get\_next\_token}]
def get_next_token(self, x, state, enc_out):
    inp_emb = self.embed_layer(x)                    # (b, emb_sz)
    h, new_state = self.decoder_cell(inp_emb, state) # h: (b, h)

    # Bilinear attention
    enc_proj = self.dense_attn(enc_out)              # (b, enc_len, h)
    h_expand = tf.expand_dims(h, axis=1)             # (b, 1, h)
    scores = tf.matmul(h_expand, enc_proj, transpose_b=True)  # (b, 1, enc_len)
    attn_weights = tf.nn.softmax(scores, axis=-1)    # (b, 1, enc_len)
    context = tf.squeeze(tf.matmul(attn_weights, enc_out), axis=1)  # (b, h)

    combined = tf.concat([h, context], axis=-1)      # (b, 2h)
    logits = self.dense(combined)                    # (b, v_sz)
    out = tf.argmax(logits, axis=-1, output_type=tf.int32)
    return out, new_state
\end{lstlisting}

\subsubsection{损失函数}

两个模型均使用稀疏 softmax 交叉熵损失函数：

\begin{equation}
    \mathcal{L} = -\frac{1}{B \cdot L} \sum_{i=1}^{B} \sum_{t=1}^{L} \log P(y_t^{(i)} \mid \mathbf{x}^{(i)}, y_{<t}^{(i)})
\end{equation}

对应代码实现：

\begin{lstlisting}[caption={损失函数}]
def compute_loss(logits, labels):
    losses = tf.nn.sparse_softmax_cross_entropy_with_logits(
            logits=logits, labels=labels)
    return tf.reduce_mean(losses)
\end{lstlisting}

%=============================================================
\section{实验结果}
%=============================================================

\subsection{任务一：普通 Seq2Seq 训练过程}

模型训练 3000 步，序列长度为 20，批量大小为 32，训练 loss 变化如下：

\begin{table}[H]
    \centering
    \caption{普通 Seq2Seq 训练 loss 变化}
    \begin{tabular}{cc}
        \toprule
        训练步数 & Loss \\
        \midrule
        0    & 3.3047 \\
        500  & 1.6266 \\
        1000 & 1.0645 \\
        1500 & 0.7858 \\
        2000 & 0.5625 \\
        2500 & 0.4624 \\
        3000（最终） & 0.3471 \\
        \bottomrule
    \end{tabular}
\end{table}

初始 loss 约为 $\ln(27) \approx 3.30$，符合随机初始化的理论预期（26个字母均匀分布）。经过 3000 步训练，loss 稳定下降至 0.347，模型收敛良好。

\subsection{任务二：Attentive Seq2Seq 训练过程}

模型训练 2000 步（少于普通模型），序列长度同为 20，训练 loss 变化如下：

\begin{table}[H]
    \centering
    \caption{Attentive Seq2Seq 训练 loss 变化}
    \begin{tabular}{cc}
        \toprule
        训练步数 & Loss \\
        \midrule
        0    & 3.2946 \\
        500  & 1.3608 \\
        1000 & 0.2376 \\
        1500 & 0.1008 \\
        2000（最终） & 0.0457 \\
        \bottomrule
    \end{tabular}
\end{table}

引入注意力机制后，模型收敛速度显著加快——在第 1000 步时 loss 已降至 0.238，而普通模型在第 3000 步才达到 0.347。

\subsection{测试结果对比}

训练完成后，对两个模型进行测试，结果如下：

\begin{table}[H]
    \centering
    \caption{两模型测试结果对比}
    \begin{tabular}{lccc}
        \toprule
        模型 & 训练步数 & 最终 Loss & 测试准确率（32样本）\\
        \midrule
        普通 Seq2Seq       & 3000 & 0.347 & 31/32 = \textbf{96.9\%}（序列长度10）\\
        Attentive Seq2Seq & 2000 & 0.046 & 28/32 = \textbf{87.5\%}（序列长度20）\\
        \bottomrule
    \end{tabular}
\end{table}

需要特别指出：两个模型的测试序列长度不同。普通 Seq2Seq 在长度为 10 的序列上测试，而 Attentive Seq2Seq 在长度为 20 的序列上测试——后者难度更高。若将测试条件统一，Attentive Seq2Seq 的实际表现预期更优。

\subsection{典型测试样本展示}

以下列出部分测试样本（普通 Seq2Seq，序列长度10）：

\begin{table}[H]
    \centering
    \caption{普通 Seq2Seq 部分测试结果}
    \begin{tabular}{ccc}
        \toprule
        输入序列 & 模型输出 & 是否正确 \\
        \midrule
        VBGMDCSMIT & TIMSCDMGBV & \checkmark \\
        GPSRQNZRSC & CSRZNQRSPG & \checkmark \\
        NYFTPYGZQI & IQZGYPTFYN & \checkmark \\
        IPQUUOTQTY & YTQTOUUQPI & \checkmark \\
        JUGPPLIWMQ & QMWILPPGUJ & $\times$ \\
        \bottomrule
    \end{tabular}
\end{table}

以下列出部分 Attentive Seq2Seq 测试样本（序列长度20）：

\begin{table}[H]
    \centering
    \caption{Attentive Seq2Seq 部分测试结果}
    \begin{tabular}{ccc}
        \toprule
        输入序列 & 模型输出 & 是否正确 \\
        \midrule
        YOMXOFXHGYVKPMHEAWGO & OGWAEHMPKVYGHXFOXMOY & \checkmark \\
        XBGJHDXBUIFCYDOVBGKE & EKGBVODYCFIUBXDHJGBX & \checkmark \\
        TODTOEJNRQKXDXIWMVYF & FYVMWIXDXKQRNJEOTDRT & \checkmark \\
        KAJIMXYUAJPVPDKNZQXA & AXQZNKDPVPJAUYXMIJAK & \checkmark \\
        BBZQHLPWCTNEQWVGJDDF & FDDJGVWQENTCWPLHQZBB & $\times$ \\
        \bottomrule
    \end{tabular}
\end{table}

%=============================================================
\section{实验分析与讨论}
%=============================================================

\subsection{注意力机制的优势分析}

\subsubsection{信息瓶颈问题的解决}

普通 Seq2Seq 将整个输入序列的信息压缩为单一固定维度向量（\texttt{enc\_state}），对于序列逆置任务，这意味着模型需要在 128 维的向量中"记住"最长 20 个字符的完整信息。当序列较长时，这一信息瓶颈导致模型难以准确记忆序列的每个位置。

注意力机制通过保留 Encoder 所有时刻的输出 $\{h_1^{\text{enc}}, \ldots, h_T^{\text{enc}}\}$，使得 Decoder 在每步解码时可以直接"查阅"输入序列的任意位置，从根本上消除了信息瓶颈。

\subsubsection{序列逆置与注意力的天然契合}

序列逆置任务要求输出序列的第 $t$ 个 token 对应输入序列的第 $(T-t+1)$ 个 token。理论上，完美的注意力模型应学会在解码第 $t$ 步时，将注意力集中于输入序列的第 $(T-t+1)$ 个位置，形成一种"反对角线"的注意力模式。这种明确的对齐关系使得注意力机制在此任务上具有天然优势。

\subsubsection{收敛速度对比}

从训练曲线可以看出，Attentive Seq2Seq 在第 1000 步时 loss 已降至 0.238，而普通 Seq2Seq 在第 3000 步时才达到 0.347。注意力模型以 \textbf{少 33\% 的训练步数}取得了\textbf{低 7 倍}的最终 loss，充分说明注意力机制对于序列对齐任务的高效性。

\subsection{关于 Teacher Forcing 与推理的差异}

训练时使用 Teacher Forcing，即 Decoder 每步的输入是真实标签，而推理时使用自回归方式，每步输入是上一步的预测结果。这种训练与推理的不一致性（exposure bias）是 Seq2Seq 模型的经典问题之一。从测试结果可以看到，模型整体上能较好地泛化，说明在本任务中该问题影响有限。

\subsection{遇到的问题与解决方案}

在实验过程中遇到了以下几个典型问题：

\begin{enumerate}
    \item \textbf{\texttt{batch\_input\_shape} 参数报错}：原代码中 \texttt{Embedding} 层使用了 \texttt{batch\_input\_shape=[None, None]} 参数，该参数在旧版 Keras（TF 1.x）中有效，但在新版 Keras 3（Python 3.12）中已被移除。解决方案：直接删除该参数，Embedding 层会自动推断输入形状。

    \item \textbf{\texttt{No gradients provided for any variable} 报错}：原代码在 \texttt{train\_one\_step} 和 \texttt{call} 方法上使用了 \texttt{@tf.function} 装饰器。在新版 Keras 3 中，\texttt{@tf.function} 将函数编译为静态图，导致 \texttt{GradientTape} 无法正确追踪 Keras 变量，所有梯度均为 \texttt{None}。解决方案：去掉相关 \texttt{@tf.function} 装饰器，以 Eager 模式运行。

    \item \textbf{Attentive Seq2Seq 推理准确率偏低}：初始实现中 \texttt{encode} 方法返回 \texttt{[enc\_out[:, -1, :], enc\_state]}，导致传入 \texttt{decoder\_cell} 的 state 列表包含两个元素，与训练时 \texttt{initial\_state=[enc\_state]} 的单元素格式不一致，造成推理阶段 state 解析错误。解决方案：修改为返回 \texttt{[enc\_state]}，使训练与推理的 state 格式保持一致。
\end{enumerate}

%=============================================================
\section{实验总结}
%=============================================================

本次实验通过字符串序列逆置这一玩具任务，完整实现并对比了普通 Seq2Seq 模型与引入双线性注意力机制的 Attentive Seq2Seq 模型。主要结论如下：

\begin{enumerate}
    \item \textbf{Seq2Seq 模型基本可行}：普通 Seq2Seq 模型在序列长度为 10 的测试集上达到 96.9\% 的准确率，证明了该架构能够有效学习序列的逆置映射关系。

    \item \textbf{注意力机制显著提升性能}：Attentive Seq2Seq 仅用 2000 步训练（少于普通模型的 3000 步）就将最终 loss 降至 0.046，相比普通模型的 0.347 低约 7 倍，在更难的长度 20 序列上仍达到 87.5\% 的准确率。

    \item \textbf{双线性注意力实现简洁有效}：通过一个可学习的线性投影层实现注意力评分，结构简单且效果良好，是 Luong Attention 的经典变体。

    \item \textbf{信息瓶颈是 Seq2Seq 的核心局限}：注意力机制通过保留并利用 Encoder 所有时刻的输出，从根本上解决了将长序列压缩为单一向量的信息瓶颈问题，是 Seq2Seq 模型最重要的改进之一，也是现代 Transformer 架构的重要前身。

    \item \textbf{工程实践中的兼容性问题}：在使用新版本深度学习框架时，需注意 API 的兼容性变化（如 \texttt{@tf.function} 与 Keras 3 的交互问题），以及训练与推理阶段 state 格式须保持一致的重要性。
\end{enumerate}

通过本次实验，深入理解了 Seq2Seq 框架的设计思想，掌握了注意力机制的数学原理与代码实现，为后续学习 Transformer 等更复杂的序列模型奠定了坚实基础。

\end{document}  